\section{The \texttt{Lens\_simple} McXtrace Component}
Simple refractive x-ray lens

\subsection*{Identification}
\begin{itemize}
  \item \textbf{Author:} Erik Knudsen
  \item \textbf{Origin:} Risoe
  \item \textbf{Date:} June 16, 2009
\end{itemize}

\subsection*{Description}
Models a stack of N refractive lenses, with a radius of curvature, r, at the apex. The model is a thin-lens approximation where photons are refracted in a the XY plane at Z=0. Absorption is generally disregarded may be handled through the use of the optional transmission parameter T, where 0\textless{}=T\textless{}=1. Thus, the lens has the focal length of f=R/(2*N*\&delta) where the x-ray refractive index is written: n = 1 - \&delta + i \&beta.

Example: Lens\_simple(xwidth=1e-5, yheight=1e-5, material\_datafile="Be.txt",N=100,r=0.3e-3)

\subsection*{Input parameters}
Parameters in \textbf{boldface} are required; the others are optional.

\begin{longtable}{p{0.22\textwidth}p{0.12\textwidth}p{0.46\textwidth}p{0.14\textwidth}}
\toprule
\textbf{Name} & \textbf{Unit} & \textbf{Description} & \textbf{Default} \\
\midrule
\endhead
xwidth & m & Width of lens aperture. & 0 \\
yheight & m & Height of lens aperture. & 0 \\
radius & m & Radius of lens aperture (overrides xwidth \& yheight). & 1e-3 \\
T & 0-1 & Transmission efficiency of the lens. & 1 \\
r & m & The radius of curvature of the lens. & 3e-4 \\
N & 1 & The number of successive lenses in the stack. & 1 \\
verbose & 0/1 & Extra information for debugging. & 0 \\
f & m & Focal length - overrides the material\_datafile - and diregards chromatic aberration. & 0 \\
material\_datafile &   & File where the material parameters for the lens may be found. Format is similar to what may be found off the NIST website. & "Be.txt" \\
\bottomrule
\end{longtable}

\subsection*{Links}
\begin{itemize}
  \item Component source code found in file \texttt{Lens\_simple.comp}.
  \item material datafile obtained from http://physics.nist.gov/cgi-bin/ffast/ffast.pl
\end{itemize}
\IfFileExists{optics/Lens_simple_static.tex}{\input{optics/Lens_simple_static.tex}}{}